\documentclass{article}

\usepackage{graphicx}
\usepackage[utf8]{inputenc}
\usepackage[T2A]{fontenc}
\usepackage[macedonian]{babel}
\usepackage{float}
\usepackage{booktabs}
\usepackage{placeins}
\usepackage{amsmath}
\usepackage{tikz}
\usepackage{pgfplots}
\usepackage{hyperref}

\pgfplotsset{compat=1.18}

\title{Компаративна анализа на управувани Kubernetes кластери од различни cloud провајдери}
\author{Стојан Делеников 259102 \and
Христина Митева 259103}
\date{Мај 2026}

\begin{document}

\maketitle

\subsection{Апстракт}

Овој проект ја анализира скалабилноста и перформансите на апликација за множење матрици поставена на две управувани Kubernetes cloud платформи, Amazon Web Services Elastic Kubernetes Service (EKS) и Microsoft Azure Kubernetes Service (AKS), конфигурирани со по три worker nodes.

Главната цел е да се испита влијанието на зголемувањето на димензиите на матриците и бројот на истовремени корисници врз латентноста, пропусниот опсег и процентот на грешки во системот. За таа цел, беа извршени load testing експерименти со користење на k6 framework, преку повеќе тест сценарија со постепено зголемување на оптоварувањето.

Во текот на експериментите беа следени метрики како време на одговор, пропусен опсег и искористеност на процесорската моќ. Добиените резултати покажуваат дека зголемувањето на комплексноста на алгоритмот и бројот на корисници има значително влијание врз перформансите и скалабилноста на системот.

Horizontal Pod Autoscaling (HPA) автоматски го зголемуваше бројот на pod инстанци при поголемо оптоварување со цел да се одржи стабилноста. 

Резултатите покажаа дека AWS EKS и Azure AKS имаат слични перформанси при помали оптоварувања, додека AWS EKS постигнува подобра латентност и пропусен опсег при потешки CPU интензивни сценарија.


\section{Вовед}

Software-as-a-Service (SaaS) апликациите сè повеќе се потпираат на технологиите за оркестрација на контејнери како Kubernetes со цел да постигнат скалабилност, доверливост, преносливост и поефикасно искористување на ресурсите. 

Управуваните Kubernetes сервиси што ги нудат cloud провајдерите овозможуваат поедноставено поставување, мониторирање и автоматско скалирање на апликациите, при што системот може динамички да се адаптира на променливи оптоварувања. 

Перформансите на контејнеризираните апликациите зависат од фактори како конфигурацијата на виртуелните машини, полисите за autoscaling, начинот на распределба на ресурсите, како и самото оптоварување.

Множењето на матрици претставува CPU интензивен проблем бидејќи неговата пресметковна комплексност расте кубно со зголемување на големината на матриците.  Поради ова, ваквиот обем на работа често се користи за тестирање и анализа на перформансите, скалабилноста и ефикасноста на cloud инфраструктури во контролирани услови.

Целта на ова истражување е да се евалуира перформансот и скалабилноста на апликацијата за множење на матрици поставен на две управувани Kubernetes платформи: AWS Elastic Kubernetes Service (EKS) и Azure Kubernetes Service (AKS). 

Анализата се фокусира на мерење на време на одговор, пропусен опсег, искористеност на процесорска моќ, како и стабилност на системот при различно оптоварување и различен број на конкурентни корисници.

Главните истражувачки прашања кои се адресираат во овој проект се:

\begin{itemize}
    \item Дали постои конзистентна разлика во латентноста и пропустната моќ помеѓу AWS EKS и Azure AKS при исти услови? 
    \item Како се менуваат перформансите со зголемување на пресметковната комплексност и бројот на конкурентни корисници? 
    \item Дали Kubernetes Horizontal Pod Autoscaling (HPA) може ефикасно да скалира при висока искористеност на процесорот?
    \item Во кој момент почнува да се намалува перформансот  за избраната конфигурација на кластерот? 


\end{itemize}
За потребите на истражувањето беше користен k6 load testing framework, преку кој беа извршени повеќе тест сценарија со постепено зголемување на оптоварувањето. Со анализа на резултатите се добива увид во однесувањето и ограничувањата на системот при зголемено оптоварување.

\section{Поврзани истражувања}

Бројни истражувања ја проучуваат скалабилноста на управувани Kubernetes Cloud околини. Horizontal Pod Autoscaling (HPA) често се користи за динамички да се прилагоди бројот на реплики на апликацијата врз однос на процесорка или мемориска искористеност. 

Едно од ваквите истражувања е трудот \textit{Container Orchestration in Multi-Cloud Environments: A Performance Evaluation} \cite{multicloud2025}, кој ја анализира Kubernetes orchestration инфраструктурата во multi-cloud околина користејќи реален e-commerce микросервисен систем. Во истражувањето се користат три cloud провајдери (вклучувајќи го GKE), при што секој кластер е конфигуриран со 5 worker nodes со исти хардверски карактеристики (4 vCPU, 16 GB RAM). Оптоварувањето се генерира со Locust и вклучува од 100 до 1000 конкурентни корисници во временски интервали од 30 минути.

Во споредба со наведеното истражување, овој проект се разликува во неколку аспекти. Меѓу главните разлики може да се издвои тоа што претходниот труд ја анализира скалабилноста на дистрибуирана микросервисна e-commerce апликација, додека овој проект се фокусира на перформансите и скалабилноста на CPU-интензивен сервис за множење на матрици. Дополнително, нивното истражување опфаќа три cloud провајдери, додека овој проект е фокусиран на споредба помеѓу AWS Elastic Kubernetes Service (EKS) и Azure Kubernetes Service (AKS). Во врска со инфраструктурата, во претходното истражување кластерите користат идентични хардверски конфигурации кај сите cloud провајдери, додека во овој проект се користат различни типови на worker јазли (general-purpose кај EKS и burstable кај AKS), што претставува ограничување при директната споредба на резултатите. Различен е и пристапот при load тестирањето, бидејќи претходното истражување користи Locust со од 100 до 1000 конкурентни корисници во интервали од 30 минути, додека тука се користи k6 framework со тест сценарија од 1 минута и различни комбинации од големина на матрици и број на конкурентни корисници. Со анализа на нивните разлики, може да се заклучи дека двете истражувања повеќе се надополнуваат отколку што можат директно да се споредуваат, бидејќи овој проект се фокусира на CPU-bound скалабилност при различни контролирани тест сценарија.

Друг релевантен труд е \textit{A Performance Evaluation of Containers Running on Managed Kubernetes Services} \cite{ferreira2019}, кој ги анализира перформансите на контејнеризирани апликации поставени на управувани Kubernetes сервиси како EKS, AKS и GKE. Ова истражувањето ги анализира метриките како процесорска моќ, меморија, мрежни ресурси, во споредба со мануелно поставена Kubernetes инфраструктура.

Наведеното истражување во својот фокус ги става однесувањето и искористеноста на системските ресурси, додека овој проект ги анализира перформансите на CPU-интензивна апликација преку анализа на латентност, пропусен опсег и autoscaling при реални тест сценарија. 

Ова го прави нивниот труд добра основа за разбирање на Kubernetes инфраструктурните перформанси, за разлика од ова истражување кое придонесува со практична анализа на скалабилност и однесување на системот при високо оптоварување.


\section{Методи}

\subsection{Методологија}

Апликацијата што се користи за тестирање е мала серверска апликација со еден endpoint \texttt{POST /compute?size=\{S\}}. За секое барање се генерираат две детерминистички матрици и се извршува нивно множење со алгоритам со три вгнездени циклуси, што создава значително процесорско оптоварување.

 Комплексноста на алгоритамот е приближно:

\[
O(n^3)
\]

каде што \(n\) ја претставува димензијата на матрицата.

\subsection{Конфигурација на кластерите}

Двата кластери користат по 3 worker nodes:
\begin{itemize}
    \item AWS EKS: \texttt{m7i-flex.large} (general-purpose класа)
    \item Azure AKS: \texttt{Standard\_B2s\_v2} (burstable класа)
\end{itemize}

Овие виртуелни машини не се еквивалентни, па затоа експериментот споредува перформанси на управувани Kubernetes платформи во реални услови. Kubernetes ресурсите се конфигурирани со \texttt{requests} од 500m CPU и 512Mi меморија, а \texttt{limits} се поставени на 1 CPU и 1Gi меморија. Horizontal Pod Autoscaler (HPA) е поставен со минимум 2 реплики, максимум 10 реплики и таргетирано процесорско оптоварување од 70\%, со цел автоматско скалирање на pod инстанците врз основа на CPU utilization.

\subsection{Тест случаи}

Load testing експериментите беа извршени со користење на k6 framework и предефинирани сценарија што комбинираат:
\begin{itemize}
    \item димензија на матрица $\in \{10, 25, 50, 100, 250, 350, 500\}$
    \item број на корисници $\in \{1, 5, 10, 25, 50, 100\}$
\end{itemize}

Вкупно се дефинирани 42 тест случаи (TC1--TC42), при што секој тест трае 1 минута, со threshold услов \texttt{p(95) < 60000 ms}.

\subsection{Метрики за евалуација}

Во анализата се следат следните метрики:
\begin{itemize}
    \item просечна латентност (avg latency)
    \item p95 латентност
    \item максимална латентност
    \item број на барања во секунда (requests/sec)
    \item вкупен број барања
    \item стапка на грешки (error rate)
\end{itemize}

Speedup метриката е дефинирана како:
\[
\text{Speedup}_{EKS/AKS} = \frac{T_{AKS}}{T_{EKS}}
\]
каде што \(T\) го претставува просечното време на одговор.

\section{Резултати}

Со изведување на експериментот, тест сценаријата успешно беа извршени до димензија на матрица со 350 x 350 елементи со оптоварување до 50 конкурентни корисници. Сценаријата со димензија од 350 x 350 елементи и 100 корисници беа прекинати поради надминување на p95 threshold условот, односно 95\% од барањата да завршат под 60 секунди. Поради тоа, сценаријата со димензија на матрица со 500 x 500 елементи не беа извршени. Добиените резултати се прикажани во следните табели и графици.

\subsection{Табели}
\begin{table}[H]
\centering
\caption{Просечна латентност според димензија на матрица}
\begin{tabular}{rrrr}
\toprule
Size & AKS avg (ms) & EKS avg (ms) & EKS better (\%) \\
\midrule
10 & 45.30 & 44.15 & 2.55 \\
25 & 47.30 & 45.81 & 3.15 \\
50 & 57.16 & 52.25 & 8.60 \\
100 & 329.64 & 169.60 & 48.55 \\
250 & 7656.41 & 7455.72 & 2.62 \\
350 & 12987.87 & 12397.19 & 4.55 \\
\bottomrule
\end{tabular}

\end{table}
\begin{table}[H]
\centering
\caption{Перформанси на матрица со димензии 350 x 350 елементи и 50 корисници}
\begin{tabular}{lrrrr}
\toprule
Platform & Avg (ms) & p95 (ms) & Error rate (\%) & Req/s \\
\midrule
AKS & 34640.36 & 59955.42 & 21.11 & 1.00 \\
EKS & 33072.33 & 59551.13 & 17.17 & 1.31 \\
\bottomrule
\end{tabular}
\end{table}
\FloatBarrier
\subsection{Графички приказ}
\begin{figure}[H]
\centering
\begin{tikzpicture}
\begin{axis}[
    width=0.9\textwidth,
    height=6cm,
    xlabel={Големина (димензија) на матрица},
    ylabel={Просечен одговор (ms)},
    legend pos=north west,
    ymajorgrids=true,
    xmajorgrids=true
]
\addplot+[mark=o, thick] coordinates {
    (10,45.3) (25,47.3) (50,57.16) (100,329.64) (250,7656.41) (350,12987.87)
};
\addlegendentry{AKS}
\addplot+[mark=square, thick] coordinates {
    (10,44.15) (25,45.81) (50,52.25) (100,169.6) (250,7455.72) (350,12397.19)
};
\addlegendentry{EKS}
\end{axis}
\end{tikzpicture}
\caption{Просечна латентност по големина (димензија) на матрица.}
\end{figure}
\begin{figure}[H]
\centering
\begin{tikzpicture}
\begin{axis}[
    width=0.9\textwidth,
    height=6cm,
    xlabel={Големина (димензија) на матрица},
    ylabel={Барања/сек},
    legend pos=north east,
    ymajorgrids=true,
    xmajorgrids=true
]
\addplot+[mark=o, thick] coordinates {
    (10,30.35) (25,30.24) (50,29.66) (100,19.69) (250,2.84) (350,0.94)
};
\addlegendentry{AKS}
\addplot+[mark=square, thick] coordinates {
    (10,30.37) (25,30.26) (50,29.83) (100,24.25) (250,2.93) (350,1.05)
};
\addlegendentry{EKS}
\end{axis}
\end{tikzpicture}
\caption{Просечна пропусносен опсег по големина (димензија) на матрица.}
\end{figure}
\begin{figure}[H]
\centering
\begin{tikzpicture}
\begin{axis}[
    width=0.9\textwidth,
    height=6cm,
    xlabel={Големина (димензија) на матрица},
    ylabel={Speedup AKS/EKS},
    ymajorgrids=true,
    xmajorgrids=true
]
\addplot+[mark=triangle, thick] coordinates {
    (10,1.026) (25,1.033) (50,1.094) (100,1.944) (250,1.027) (350,1.048)
};
\addplot[dashed] coordinates {(10,1.0) (350,1.0)};
\end{axis}
\end{tikzpicture}

\caption{Speedup ($T_{AKS}/T_{EKS}$), 
каде вредности поголеми од 1 укажуваат на подобри перформанси на EKS.}
\end{figure}
\FloatBarrier

\section{Дискусија}
Резултатите јасно покажуваат дека со зголемување на големината на матрицата и бројот на конкурентни корисници, латентноста нагло расте, што е очекувано за алгоритам со комплексност $O(n^3)$.  

Кај помали матрици со димензии од 10 x 10 до 50 x 50 елементи, двете платформи имаат слични перформанси, со некои мали разлики. Кај поголеми матрици, особено матрица со димензии од 100 x 100 елементи, EKS покажува значително пониска просечна латентност во споредба со AKS. Кај најголемите матрици, матрици со димензии од 250 x 250 и 350 x 350 елементи, двете платформи го достигнуваат максималниот капацитет при што се забележува опаѓање на пропусноста и раст на латентноста.

Дополнително, бројот на реплики не беше фиксен во текот на експериментите. При сценаријата со помали матрици и помал број конкурентни корисници, апликацијата најчесто работеше со две реплики. Со зголемување на големината на матриците и интензитетот на оптоварувањето, Horizontal Pod Autoscaler (HPA) постепено го зголемуваше бројот на реплики, достигнувајќи до десет реплики кај најоптоварените сценарија. Покрај ова автоматско скалирање, кај најголемите матрици и највисокиот број корисници се забележуваше раст на латентноста, што укажува дека дополнителните реплики не можеа целосно да го компензираат зголемувањето на оптоварувањето.

Со зголемување на оптоварувањето, p95 латентноста значително расте, при што кај најтешките извршени сценарија со матрица од 350 x 350 елементи и 50 корисници се приближува до прагот од 60 секунди. Истовремено, се забележува и зголемување на стапката на грешки, особено при повисок број истовремени корисници. Во оваа зона AWS EKS сепак задржува мала предност со што покажа подобра стабилност и нешто пониска латентност. Како можни причини за ова се разликите во процесорското однесување на виртуелните машини од различни класи, како и доцнењето на HPA механизмот при ненадејно зголемување на оптоварувањето.

Клучна забелешка е дека класите на виртуелни машини по дизајн се различни, односно general-purpose на EKS наспроти burstable на AKS. Ова може суштински да влијае врз процесорските перформанси, динамиката на ограничување на процесорот (throttling) и латентноста кај најбавните барања при високо оптоварување.

Според наведеното, дел од забележаните разлики веројатно произлегуваат од класата на виртуелните машини, а не само од однесување на управуваниот Kubernetes control plane node.

\section{Заклучок}

Ова истражување направи споредба помеѓу AWS EKS и Azure AKS при извршување на CPU-интензивен алгоритам за множење на матрици. Тестирањето беше извршено со k6 сценарија, при различен број на конкурентни корисници и различни големини на матрици, со следење на клучни метрики како латентност, пропусен опсег и стапка на грешки.

Резултатите покажаа дека при мал товар двете платформи имаат слично и стабилно однесување, додека при зголемен товар перформансите се намалуваат, односно латентноста расте, пропусноста опаѓа, а стапката на грешки се зголемува. Во потешките сценарија, AWS EKS покажа мала предност во однос на Azure AKS, особено при високo процесорско оптоварување и услови блиски до максималниот капацитет на системот.

Тестовите покажаа дека скалирањето помага да се задржат подобри перформанси, односно помала латентност и стабилен пропусен опсег. Дополнително, управуваните Kubernetes платформи овозможуваат лесно поставување и скалирање на апликации. Од експериментот може да се заклучи и дека general-purpose виртуелните машини се подобар избор за worker nodes во Kubernetes околина.

Главно ограничување на извршената анализа е тоа што не беа користени целосно исти виртуелни машини на двете платформи, па дел од разликите може да се должат на инфраструктурните разлики.

Од добиените резултати од овој експеримент може да се заклучи дека AWS EKS се покажа како подобар избор кога приоритет се стабилни перформанси и пониска латентност при висок товар.


\begin{thebibliography}{9}
\bibitem{multicloud2025}
M. K. Pasupuleti,
\textit{Container Orchestration in Multi-Cloud Environments: A Performance Evaluation},
\textit{International Journal of Academic and Industrial Research Innovations (IJAIRI)}, vol. 05, pp. 327--340, 2025. DOI: \href{https://doi.org/10.62311/nesx/rphcrcscrcec2}{10.62311/nesx/rphcrcscrcec2}.

\bibitem{ferreira2019}
A. Pereira Ferreira and R. Sinnott,
\textit{A Performance Evaluation of Containers Running on Managed Kubernetes Services},
In: \textit{2019 IEEE International Conference on Cloud Computing Technology and Science (CloudCom)}, pp. 199--208, 2019. DOI: \href{https://doi.org/10.1109/CloudCom.2019.00038}{10.1109/CloudCom.2019.00038}.
\end{thebibliography}

\end{document}
